\documentclass[11pt]{article}

\usepackage[T1]{fontenc}
\usepackage[margin=1in]{geometry}
\usepackage{amsmath}
\usepackage{amssymb}
\usepackage{physics}
\newcommand{\kket}[1]{|#1\rangle\!\rangle}
\usepackage{bbold}
\usepackage{booktabs}
\usepackage[hidelinks]{hyperref}

\title{QuDPy-FDGF Solver API and Implementation Reference}
\author{Mathieu Desmarais, Sim\'on Paiva-Ortega, Hao Li, Eric R. Bittner, Carlos Silva-Acu\~{n}a}
\date{Software version 1.0.0}

\begin{document}

\sloppy

\maketitle

\begin{abstract}
This document collects the software-facing reference material for the QuDPy-FDGF solver. It describes the model contracts, public solver operations, observable specifications, pathway objects, numerical backends, and plotting utilities. The scientific manuscript contains the mathematical formulation, validation benchmarks, and physical examples; this companion document records API details that may evolve with the software.
\end{abstract}

\tableofcontents
\newpage

\section*{Notation and conventions}
The notation follows the companion manuscript. Operators on the Hilbert space carry hats ($\hat{H}$, $\hat{\rho}$, $\hat\mu$, $\hat L_c$, $\hat{O}$); calligraphic symbols denote superoperators ($\mathcal{L}$, $\mathcal{U}$, $\mathcal{R}_\eta$, $\mathcal{V}$, $\mathcal{D}_c$), and the same symbol denotes a superoperator and its matrix acting on the column-stacked vector $\operatorname{vec}(\hat{\rho})=\kket{\hat{\rho}}$; $\mathbb{1}$ is the identity. Dissipative channels are labeled by $c$, with jump operator $\hat L_c$, dissipator $\mathcal{D}_c[\hat\rho]=\hat L_c\hat\rho\hat L_c^\dagger-\frac{1}{2}\{\hat L_c^\dagger\hat L_c,\hat\rho\}$, and rate $\gamma_c$ kept separate: $\mathcal{L}[\hat\rho]=-i[\hat H,\hat\rho]+\sum_c\gamma_c\mathcal{D}_c[\hat\rho]$. The rate $\gamma_c$ is the Lindblad coefficient, not necessarily a measured rate: a pure-dephasing rate $\gamma_\phi=1/T_2^*$ requires $\gamma_c=\gamma_\phi/2$ for $\hat L_c=\hat\sigma_z$ and $\gamma_c=2\gamma_\phi$ for a projector $\ket{e}\bra{e}$ (Sec.~\ref{sec:quick-start}). The raising and lowering parts of the light--matter operator are $\hat\mu_\pm$; in the software interface they are supplied as the transition blocks $J_\pm$ (\texttt{j\_plus\_array}, \texttt{j\_minus\_array}). The pathway prefactor is $c_p$ and the detection efficiency $\varepsilon_{\mathrm{det}}$; $\eta$ is reserved for the resolvent regularization. Energies and frequencies share one unit, $\hbar=1$ (eV in the examples).


\section{Quick start}
\label{sec:quick-start}

\subsection{Installation}
QuDPy-FDGF requires Python~$\geq 3.10$. The distribution is named \texttt{qudpy-fdgf} and the import name is \texttt{qudpy\_fdgf}. From a clone of the repository, install it in editable mode:
\begin{verbatim}
git clone https://github.com/SilvaScience/QuDPy-FDGF.git
cd QuDPy-FDGF
pip install -e .              # core: NumPy, SciPy, Matplotlib
pip install -e ".[examples]"  # optional: QuTiP, UFSS, IPython (notebooks)
pip install -e ".[dev]"       # optional: pytest
\end{verbatim}
The core dependencies are NumPy~$\geq 1.24$, SciPy~$\geq 1.10$, and Matplotlib~$\geq 3.7$; Matplotlib is used only by \texttt{SpectroscopyPlotter}. QuTiP is used by the example notebooks to build operators, and \texttt{ufss.DiagramGenerator} generates the fifth-order pathways of the chromophore example. The optional extra \texttt{[manybody]} installs TeNPy for the experimental modules, which are not part of this release. The reference output below was obtained with Python~3.13.14, NumPy~2.5.0, SciPy~1.18.0, and Matplotlib~3.11.0.

\subsection{A complete calculation}
The script below treats an open two-level system with transition energy $2~\mathrm{eV}$, radiative decay at rate $\gamma_1$, and pure dephasing at rate $\gamma_\phi$. It performs the five steps common to every calculation: build the matrices, adapt them to the model contract, feed the model to the solver, declare the pathways and the protocol, and generate the spectrum.
% Distributed as examples/quickstart.py.
\begin{verbatim}
import numpy as np
from qudpy_fdgf import EigenbasisKModel, FrequencyPathway, SpectroscopySolver
from qudpy_fdgf.protocols import (PropagationInterval, SpectroscopyProtocol,
                                  standard_nq_protocol)

# 1. Matrices (hbar = 1, eV), basis ordering (|e>, |g>)
sigma_z = np.diag([1.0, -1.0]).astype(complex)
sigma_x = np.array([[0, 1], [1, 0]], dtype=complex)
sigma_m = np.array([[0, 0], [1, 0]], dtype=complex)   # |g><e|
H, mu = 1.0 * sigma_z, 0.5 * sigma_x
gamma_1, gamma_phi = 0.010, 0.005                     # 1/T1, 1/T2* (eV)

# 2. Model contract: (operator, rate) pairs, rate kept out of the operator
model = EigenbasisKModel(H, mu, c_ops_raw=(
    (sigma_m, gamma_1),          # population decay at gamma_1
    (sigma_z, gamma_phi / 2),    # sigma_z dephases coherences at 2x its rate
))

# 3. Solver
solver = SpectroscopySolver(backend="dense", eta=1e-4)
solver.feed_model(model)

# 4-5. Linear response: one ket interaction, one frequency interval
linear = FrequencyPathway("linear", interactions=[{"label": "Ku"}],
                          component="linear")
protocol = SpectroscopyProtocol(
    intervals=[PropagationInterval("omega1", "frequency")])
w = np.linspace(1.9, 2.1, 4001)
res = solver.generate_spectrum(protocol, {"omega1": w}, pathways=[linear])
A = res.pathways["linear"].imag      # absorptive part
half = w[A >= A.max() / 2]
Gamma_2 = gamma_1 / 2 + gamma_phi
print(f"linear peak      : {w[np.argmax(A)]:.4f} eV")
print(f"FWHM (numerical) : {half[-1] - half[0]:.4f} eV")
print(f"2(Gamma_2 + eta) : {2 * (Gamma_2 + solver.eta):.4f} eV")

# Third-order rephasing (R1 + R2) at t2 = 0: omega_1 and omega_3 resolved
R1 = FrequencyPathway("R1", interactions=[
    {"label": "Bu", "pulse_index": 0}, {"label": "Ku", "pulse_index": 1},
    {"label": "Bd", "pulse_index": 2}], component="rephasing")
R2 = FrequencyPathway("R2", interactions=[
    {"label": "Bu", "pulse_index": 0}, {"label": "Bd", "pulse_index": 1},
    {"label": "Ku", "pulse_index": 2}], component="rephasing")
protocol_2d = standard_nq_protocol(order=1, nq_interval=1, detection_interval=3,
                                   n_interactions=3, nq_axis="omega1q")
axes = {"omega1q": np.linspace(-2.2, -1.8, 81),
        "omega3": np.linspace(1.8, 2.2, 81)}
res2 = solver.generate_spectrum(protocol_2d, axes, pathways=[R1, R2],
                                fixed_coordinates={"t2": 0.0})
S2 = res2.components["rephasing"]
i, j = np.unravel_index(np.argmax(np.abs(S2)), S2.shape)
print(f"rephasing peak   : (omega1q, omega3) = "
      f"({res2.axis_values[0][i]:.3f}, {res2.axis_values[1][j]:.3f}) eV")
print(f"|S| at peak      : {np.abs(S2[i, j]):.2f}")
print(f"2 mu^4/(Gamma_2 + eta)^2 : {2 * 0.5**4 / (Gamma_2 + solver.eta)**2:.2f}")
\end{verbatim}

\subsection{Expected output}
The script runs in about two seconds on a laptop and prints
\begin{verbatim}
linear peak      : 2.0000 eV
FWHM (numerical) : 0.0202 eV
2(Gamma_2 + eta) : 0.0202 eV
rephasing peak   : (omega1q, omega3) = (-2.000, 2.000) eV
|S| at peak      : 1225.37
2 mu^4/(Gamma_2 + eta)^2 : 1225.37
\end{verbatim}
Each numerical line is followed by its analytical value. The absorptive line has a full width at half maximum $2(\Gamma_2+\eta)$, with $\Gamma_2=\gamma_1/2+\gamma_\phi$; the width confirms that each supplied rate $\gamma_c$ multiplies $\mathcal{D}_c[\hat\rho]$ unchanged, so that the $\hat\sigma_z$ channel must carry $\gamma_\phi/2$ (a diagonal channel $\hat L_c=\sum_x l_x\ket{x}\bra{x}$ damps $\ket{x}\bra{y}$ at the rate $\gamma_c(l_x-l_y)^2/2$). The rephasing signal peaks on the antidiagonal quadrant $(\omega_{1q},\omega_3)=(-2,2)~\mathrm{eV}$, where each of the two pathways contributes $|\hat\mu_{eg}|^4/(\Gamma_2+\eta)^2$. Replacing \texttt{backend="dense"} by \texttt{backend="sparse\_sector"} gives the same output.

\subsection{Next steps and tests}
\begin{itemize}
\item \texttt{examples/quickstart.py}: the listing above.
\item \texttt{examples/example1\_two\_level\_open.ipynb}: Example~1 of the manuscript, with populations, action detection, and integrated jump observables (Sec.~\ref{sec:observables}); the listing of the manuscript is extracted from its cells tagged \texttt{listing}.
\item \texttt{examples/example2\_xxz\_chain\_thermal.ipynb} and \texttt{examples/example3\_chi5\_2q.ipynb}: Examples~2 and~3.
\item \texttt{python -m pytest}: runs the test suite in \texttt{tests/}, including the analytical checks of the manuscript (about 17~s with a single BLAS thread, \texttt{OMP\_NUM\_THREADS=1}); \texttt{tests/api/} checks the contracts and options of the library, \texttt{tests/physics/} the closed-form results and limiting cases, and \texttt{tests/regression/} compares eight small calculations with stored outputs, which any change of the numerical core must reproduce.
\item \texttt{validation/}: analysis notebooks (convergence, analytical limits, backend agreement), benchmark notebooks and scripts (time, memory, comparison with the original QuDPy), the scripts that redraw the figures of the manuscript, the shared models, and the stored results (see \texttt{validation/README.md}).
\end{itemize}
Section~\ref{app1} documents the model contract, the solver, and the backends; Sec.~\ref{sec:pathway-conventions} fixes the sign and ordering conventions of the pathways.


\section{Solver API and implementation reference}
\label{app1}
\subsection{\texttt{SectorModel}: external model contract}

\noindent \texttt{SectorModel} is the block-based interface between an external physical model and \texttt{SpectroscopySolver}. The model defines the physical truncation, basis, operators, initial state, and dissipative channels; the solver validates, assembles, and propagates these objects. The truncated Hilbert space is a direct sum,
\begin{equation}
    \mathcal{H}=\bigoplus_{s\in\mathcal{S}}\mathcal{H}_s,
    \qquad \dim(\mathcal{H}_s)=d_s,
\end{equation}
and a block mapping source sector $s$ to target sector $t$ has the form
\begin{equation}
    \hat{A}_{ts}=\hat{P}_t \hat{A} \hat{P}_s\in\mathbb{C}^{d_t\times d_s}.
\end{equation}
Sector labels are opaque hashable objects and sector dimensions need not be equal.

\subsubsection{Required structural methods}

The minimal contract used by \texttt{feed\_model} is
\begin{itemize}
    \item \texttt{sectors()} and \texttt{dimension(sector)};
    \item \texttt{hamiltonian\_blocks(source)}, returning a mapping from target sectors to Hamiltonian blocks;
    \item \texttt{transition\_blocks(...)}, with source sector, operator name, and direction (\texttt{"plus"} or \texttt{"minus"});
    \item \texttt{observable\_blocks(...)}, returning the blocks of a named detection operator for a source sector.
\end{itemize}
Hamiltonian blocks may be diagonal or inter-sector. Operator blocks may be dense arrays, sparse matrices, or matrix-free linear operators exposing their shape and forward/adjoint actions.

\subsubsection{Initial and equilibrium states}

The model method \texttt{initial\_condition(context=None)} returns either a \texttt{PureState} or a \texttt{DensityState}. A pure state stores one sector label and a vector in that sector. A density state stores blocks under keys \texttt{(ket\_sector, bra\_sector)}:
\begin{equation}
    \rho_{ab}=\hat{P}_a\hat{\rho} \hat{P}_b\in\mathbb{C}^{d_a\times d_b}.
\end{equation}
Both populations ($a=b$) and cross-sector coherences ($a\neq b$) are supported. A block may be supplied as a dense matrix through \texttt{DenseDensityBlock}, or in factored form through \texttt{LowRankDensityBlock(X,Y,S)}, representing
\begin{equation}
    \rho_{ab}=X S Y^\dagger,
\end{equation}
with $S$ optional. During backend construction, the assembled state is normalized and checked for compatible dimensions, Hermiticity, unit trace, and positive semidefiniteness within numerical tolerance.

Thermodynamic input is normalized to
\begin{verbatim}
ThermodynamicContext(
    temperature,
    ensemble="canonical",
    chemical_potentials={},
    k_ensemble="independent",
    metadata={},
)
\end{verbatim}
The supported ensemble labels are canonical and grand canonical; the $k$ ensemble is independent or global. A missing context does not imply zero temperature or equilibrium: it delegates the default state to the model. When a context is supplied, the optional equilibrium-state method remains model-side; the solver does not infer a thermal state from the Hamiltonian. The provided adapters build the Gibbs state of $\hat H$ without reference to the collapse channels; whether it is stationary under the generator is checked by the solver (Sec.~\ref{sec:stationarity}).

\subsubsection{Transition, detection, and collapse operators}

For a named interaction operator, the model supplies the physical decomposition
\begin{equation}
    \hat{\mu}=\hat{\mu}_+ + \hat{\mu}_-.
\end{equation}
The solver does not infer this decomposition from energy differences or sector labels. The separate \texttt{observable\_blocks} method exposes named detection operators, so propagation and readout operators need not coincide.

The backends check the operators they receive. When they are built, $\hat H x$ and $\hat H^\dagger x$ are compared for two random vectors, and a non-Hermitian Hamiltonian raises \texttt{ModelContractError} (backend option \texttt{hermiticity\_tolerance}, default $10^{-10}$, relative; the check is skipped for matrix-free blocks that provide no adjoint action). At the first use of each interaction operator, $\hat\mu_-x$ is compared with $\hat\mu_+^\dagger x$, and a mismatch emits a \texttt{ModelConsistencyWarning} (option \texttt{transition\_tolerance}, default $10^{-8}$), because $\hat\mu=\hat\mu_++\hat\mu_-$ is then not Hermitian and the detected signal is not that of a physical observable. The adapters supplied with the package satisfy both conditions by construction.

The optional collapse-channel method returns \texttt{CollapseChannel} objects. A channel stores a name, a non-negative rate $\gamma_c$, and blockwise $\hat{L}_c$, with the convention
\begin{equation}
    \mathcal{L}[\hat{\rho}]=-i[\hat{H},\hat{\rho}]+\sum_c\gamma_c\,\mathcal{D}_c[\hat{\rho}],
    \qquad
    \mathcal{D}_c[\hat{\rho}]=\hat{L}_c\hat{\rho} \hat{L}_c^\dagger-\frac{1}{2}
    \{\hat{L}_c^\dagger \hat{L}_c,\hat{\rho}\}.
\end{equation}
The rate must therefore not be folded into the supplied operator blocks.

\subsubsection{Requirements and capabilities}

The optional \texttt{requirements()} and \texttt{capabilities()} methods advertise the requested state kind, generator kind, propagation domains, exactness, matrix-free constraint, and related metadata. The solver combines this declaration with the effective initial state and collapse channels before selecting or validating a backend. Thus \texttt{SectorModel} is a data contract, not a prescribed physical model.





\subsection{\texttt{EigenbasisKModel}: provided $k$-space adapter}

\noindent \texttt{EigenbasisKModel} constructs a \texttt{SectorModel}-compatible object from site-basis Hamiltonians and operators. A single $(d,d)$ Hamiltonian produces one sector; an $(N_k,d,d)$ stack produces independent sectors $\{k_0,\ldots,k_{N_k-1}\}$. This adapter introduces no inter-$k$ coupling.

\subsubsection{Constructor inputs}
The constructor has the form 
\begin{verbatim}
EigenbasisKModel(
    H_model,
    interaction_op_array,
    *,
    c_ops_raw=(),
    detection_op_array=None,
    observable_op_arrays=None,
    j_plus_array=None,
    j_minus_array=None,
    rwa_tol=1e-6,
    transition_window=None,
    k_weights=None,
    boltzmann_constant=8.617333262e-5,
)
\end{verbatim}

\subsubsection{Diagonalization and basis transformation}
For each k-point, the adapter diagonalizes the site-basis Hamiltonian,

\begin{equation}
    \hat{H}_{site}^{(k)}U^{(k)} = U^{(k)} E^{(k)}. 
\end{equation}

where $E^{(k)}$ is diagonal. Every supplied interaction, detection, observable, and collapse operator is transformed according to
\begin{equation}
    \hat{A}_{\mathrm{eig}}^{(k)} = \big(U^{(k)}\big)^\dagger \hat{A}_{\mathrm{site}}^{(k)} U^{(k)}.
\end{equation}
These transformations are performed once during model construction.

\subsubsection{Automatic transition-operator decomposition}
\label{sec:auto_decomposition}
When explicit \texttt{j\_plus\_array} and \texttt{j\_minus\_array} inputs are not provided, the transformed interaction operator is separated according to the eigenenergy differences. 

For a matrix element connecting source eigenstate j to target eigenstate i, define
\begin{equation}
    \Delta E_{ij}^{(k)} = E_i^{(k)} - E_j^{(k)}
\end{equation}
The automatic decomposition is

\begin{align}
\left(\hat{\mu}_{+}^{(k)}\right)_{ij}
&=\begin{cases}\left(\hat{\mu}^{(k)}\right)_{ij},
&\Delta E_{ij}^{(k)}>\texttt{rwa\_tol},\\0,
&
\text{otherwise},
\end{cases}
\\[4pt]
\left(\hat{\mu}_{-}^{(k)}\right)_{ij}
&=
\begin{cases}
\left(\hat{\mu}^{(k)}\right)_{ij},
&
\Delta E_{ij}^{(k)}<-\texttt{rwa\_tol},\\
0,&\text{otherwise}.
\end{cases}
\end{align}

Matrix elements satisfying
\begin{equation}
    |\Delta E_{ij}^{(k)}| \leq \text{rwa\_tol}
\end{equation}
are omitted from both components; this includes the static dipole moments and the elements between degenerate states. Models requiring explicit manifold assignments, degeneracy handling, or external pathway conventions may instead provide both \texttt{j\_plus\_array} and \texttt{j\_minus\_array}. The reported decomposition is then \texttt{"automatic\_energy"} or \texttt{"explicit"}.

\paragraph{Bandwidth window and report.}
In the impulsive limit every transition of positive energy is treated as resonant with the field, however small its Bohr frequency, which is wrong for a low-frequency transition (vibrational, or intra-manifold) of a model driven by an optical pulse. The argument \texttt{transition\_window=}$(\omega_{\min},\omega_{\max})$ keeps only the elements with $\omega_{\min}\leq|\Delta E_{ij}|\leq\omega_{\max}$, in the units of the Hamiltonian: the band that the pulses of the experiment can drive. It requires $0\leq\omega_{\min}<\omega_{\max}$ and cannot be combined with explicit \texttt{j\_plus\_array} and \texttt{j\_minus\_array}. The method \texttt{transition\_summary()} returns a dictionary with the fractions of $\sum_{ij}|\mu_{ij}|^2$ that are kept (\texttt{kept\_fraction}), static (\texttt{static\_fraction}: diagonal elements of non-degenerate states), inside degenerate subspaces (\texttt{degenerate\_fraction}) and excluded by the window (\texttt{outside\_window\_fraction}), together with \texttt{n\_degenerate\_states}; the degenerate weight is the Frobenius weight of the zero-frequency blocks, which does not depend on the basis chosen inside a degenerate subspace. If the degenerate weight exceeds $10^{-6}$ of the total, the constructor emits a \texttt{ModelConsistencyWarning}: the field does not couple these states, and the explicit arrays must be used if it should. \texttt{ExcitationSectorModel.transition\_summary()} reports \texttt{"explicit\_sector"}.

\subsubsection{Detection and named observables}

For every sector $k$, the adapter returns only intra-sector blocks:
\begin{align}
\texttt{hamiltonian\_blocks(k)}
&=
\{\texttt{k}:\hat{H}_{\mathrm{eig}}^{(k)}\},\\
\texttt{transition\_blocks(\text{"light\_matter"},\text{"plus"},k)}
&=
\{\texttt{k}:\hat{\mu}_{+}^{(k)}\},\\
\texttt{transition\_blocks(\text{"light\_matter"},\text{"minus"},k)}
&=
\{\texttt{k}:\hat{\mu}_{-}^{(k)}\}.
\end{align}

The default detection operator is the transformed \texttt{detection\_op\_array}; if it is omitted, the interaction operator is reused. The observable mapping adds independently named operators, transformed in the same way. Their names are returned by \texttt{observable\_names()}. When named observables are configured, requesting an unknown name raises an error; without that mapping, arbitrary names retain the legacy fallback to the single detection operator.


\subsubsection{Initial and equilibrium states}

The default initial condition is a mixed state distributed over the $k$ sectors according to normalized \texttt{k\_weights}. Within each sector, the population occupies the lowest eigenstate:

\begin{equation}
    \rho^{(0)}_{kk}=w_k\ket{0_k}\bra{0_k},
    \qquad \sum_k w_k=1.
\end{equation}
For a supplied thermodynamic context, the equilibrium state is constructed from the eigenergies of each sector. At positive temperature,
\begin{equation}
    \rho_{kk}^{\mathrm{eq}}=w_k\sum_i
    \frac{\exp[-\beta E_i^{(k)}]}{Z_k}\ket{i_k}\bra{i_k},
    \qquad Z_k=\sum_i\exp[-\beta E_i^{(k)}],
\end{equation}
with 
\begin{equation}
    \beta = \frac{1}{k_B T}
\end{equation}

where the configurable \texttt{boltzmann\_constant} must use units consistent with the Hamiltonian. At zero temperature, degenerate lowest states within each sector share its weight. The present adapter treats the $k$ sectors independently and preserves the fixed $k$ weights.

\subsubsection{Dissipative channels and backend requirements}
Each raw collapse-operator pair ($\hat{L}_c$, $\gamma_c$) is transformed into the eigenbasis of every $k$-sector. The resulting channel has a generated name such as \texttt{c\_op\_0} and contains only intra-sector blocks, 
\begin{equation}
    \hat{L}_{c,\mathrm{eig}}^{(k)} = \big( U^{(k)} \big)^\dagger \hat{L}_{c,\mathrm{site}}^{(k)} U^{(k)}.
\end{equation}
If at least one collapse channel is supplied, the adapter advertises a Lindblad generator; otherwise it advertises unitary evolution. Its default state representation is mixed because the independent $k$ sectors are combined incoherently.

\subsection{\texttt{ExcitationSectorModel}: explicit manifold adapter}

\noindent \texttt{ExcitationSectorModel} is the adapter for models already organized into excitation or other user-defined sectors. Unlike \texttt{EigenbasisKModel}, sector dimensions may differ and rectangular inter-sector operator blocks are supported explicitly.

\subsubsection{Constructor and block convention}

\begin{verbatim}
ExcitationSectorModel(
    hamiltonian_blocks,
    raising_blocks,
    *,
    lowering_blocks=None,
    detection_blocks=None,
    observable_op_blocks=None,
    c_ops_raw=(),
    initial_sector=None,
    initial_state_index=0,
    boltzmann_constant=8.617333262e-5,
)
\end{verbatim}
The Hamiltonian mapping contains one dense Hermitian block per sector. All operator mappings use \texttt{(target, source)} sector keys. Each sector Hamiltonian is diagonalized as
\begin{equation}
    \hat{H}_s V_s=V_s E_s,
\end{equation}
and every rectangular block is transformed consistently,
\begin{equation}
    \hat{A}_{ts}^{\mathrm{eig}}=V_t^\dagger \hat{A}_{ts}V_s.
\end{equation}

If \texttt{lowering\_blocks} is omitted, it is generated as the blockwise adjoint of \texttt{raising\_blocks}. If \texttt{detection\_blocks} is omitted, the sum of raising and lowering blocks is used. The mapping \texttt{observable\_op\_blocks} exposes additional named observables through the same block convention. Because the transition blocks are supplied directly, \texttt{transition\_decomposition()} reports \texttt{"explicit\_sector"}.

\subsubsection{Initial and thermal states}

The default initial sector is the one containing the global lowest eigenenergy unless \texttt{initial\_sector} is specified; \texttt{initial\_state\_index} selects the eigenstate inside that sector. For a thermodynamic context, the adapter constructs one global canonical distribution over all retained sector eigenstates,
\begin{equation}
    \hat{\rho}^{\mathrm{eq}}=
    \frac{1}{Z}\sum_{s,i}\exp(-\beta E_{si})
    \ket{s,i}\bra{s,i},
    \qquad Z=\sum_{s,i}\exp(-\beta E_{si}).
\end{equation}
At zero temperature, the population is shared among all globally degenerate ground states. The configurable Boltzmann constant again fixes the energy--temperature unit conversion.








\subsection{\texttt{SpectroscopySolver}: numerical orchestrator}

\texttt{SpectroscopySolver} is the public orchestrator of the solver. It connects an external SectorModel to a numerical backend, manages the active pathways responses, and assembles multidimensional spectra. The solver does not contain model-specific physics. The physical information is supplied by the SectorModel, while the numerical propagation is delegated to the selected backend.

\subsubsection{Constructor}

\begin{verbatim}
SpectroscopySolver(
    backend="sparse_sector",
    eta=0.05,
    requirements=None,
    thermodynamic_context=None,
    check_stationarity="warn",
    stationarity_tolerance=1e-8,
    **backend_options
)
\end{verbatim}
The arguments \texttt{check\_stationarity} and \texttt{stationarity\_tolerance} control the check of the reference state described below.

When \texttt{backend="auto"} is selected, the solver compares the effective model requirements with the declared capabilities of the registered backends. Built-in and custom backends are listed by \texttt{available\_backends()} and custom implementations are added with \texttt{register\_backend(name, backend\_class)}.


\subsubsection{Model registration and backend construction}


The method \texttt{feed\_model(model, context=None)} connects a \texttt{SectorModel} to the solver. The method performs the following operations:
\begin{enumerate}
    \item validates the structural model contract;
    \item resolves the thermodynamic context, initial state, and collapse channels;
    \item updates the effective state and generator requirements;
    \item selects or validates a compatible numerical backend;
    \item constructs the sector layout, generator, and operator blocks;
    \item checks that the reference state is stationary (Sec.~\ref{sec:stationarity}).
\end{enumerate}
The optional \texttt{context} argument overrides the thermodynamic context provided to the constructor when it is not \texttt{None}. The method return the solver instance itself, allowing chained construction:
\begin{verbatim}
    solver = (
    SpectroscopySolver(backend="sparse_sector")
    .feed_model(model, context=context))
\end{verbatim}

Pathway calculations cannot be performed before \texttt{feed\_model} has successfully constructed the backend.

\subsubsection{Reference-state check}
\label{sec:stationarity}

The perturbative response of the manuscript starts from a reference state $\hat\rho_{\mathrm{ref}}$ that does not evolve before the first pulse, $\mathcal{L}\hat\rho_{\mathrm{ref}}=0$. The model supplies this state and nothing in the Hamiltonian guarantees the property: the Gibbs state of $\hat H$ is stationary for a closed model, but with collapse channels it is stationary only if the channels satisfy the detailed balance of $\hat H$ at that temperature. At the end of \texttt{feed\_model} the solver therefore applies the generator once to the reference state and computes the relative residual
\begin{equation}
    r=\frac{\lVert\mathcal{L}\kket{\hat\rho_{\mathrm{ref}}}\rVert}{\lVert\kket{\hat\rho_{\mathrm{ref}}}\rVert},
\end{equation}
with Frobenius norms and in the units of the Hamiltonian and of the rates. The constructor argument \texttt{check\_stationarity} selects the behavior when $r$ exceeds \texttt{stationarity\_tolerance} (default $10^{-8}$): \texttt{"warn"} (default) emits a \texttt{StationarityWarning}, \texttt{"error"} raises \texttt{StationarityError}, and \texttt{"off"} skips the check. The method \texttt{stationarity\_residual()} returns $r$ at any time, and \texttt{summary()} reports it. For example, an open two-level system at $10~\mathrm{K}$ with the channel $\hat L=\ket{g}\bra{e}$ has a Gibbs state with $9\%$ of excited population that the channel empties: $r=0.03$, and the linear response at resonance is $9\%$ below that of the stationary state, which is the ground state. The check does not construct a stationary state; a state that fails it must be replaced by the user.



\subsubsection{Pathway registration}

The active pathways are registered with \texttt{set\_pathways(pathways)}. The input may contain \texttt{FrequencyPathway} objects or mappings from which these objects can be constructed. Names must be unique and at least one pathway must be supplied. The method returns the solver instance. Active pathways are queried with \texttt{get\_pathways(component=None)}.
\newline

If \texttt{component} is omitted, all active pathways are returned. If a component is provided, only pathways assigned to that component are returned. The aliases \texttt{"nonrephasing"} and \texttt{"nonrephase"} are normalized to the internal component name \texttt{"unrephasing"}. 
\newline

The method \texttt{pathway\_summary()} returns serializable pathway metadata without exposing the internal model or backend objects.

\subsubsection{Propagation protocols}

A \texttt{SpectroscopyProtocol} is an ordered tuple of \texttt{PropagationInterval} objects. Each interval has a unique name, a domain in \texttt{"time"}, \texttt{"frequency"}, or \texttt{"identity"}, an optional coherence order, and an optional interval-specific broadening \texttt{eta}. The protocol validates pathway length, coherence-order magnitude, and coordinates. It separately exposes all coordinate names and its frequency- and time-axis names.

The helper
\begin{verbatim}
standard_nq_protocol(
    order,
    nq_interval,
    detection_interval,
    *,
    n_interactions=3,
    nq_axis=None,
    detection_axis="omega3",
)
\end{verbatim}
constructs the common two-frequency-axis protocol, with the remaining propagation intervals kept in the time domain.

\subsubsection{Single-pathway calculation}

The method \texttt{calc\_pathway(pathway, protocol, coordinates)} evaluates one pathway at one coordinate point. 

The input are:
\begin{itemize}
    \item \texttt{pathway}: a \texttt{FrequencyPathway} object or the name of an active pathway;
    \item \texttt{protocol}: a SpectroscopyProtocol describing the propagation intervals;
    \item \texttt{coordinates}: a mapping from interval names to numerical time or frequency values. 
\end{itemize}

The protocol determines which intervals are propagated in the time or frequency domain. Every non-identity interval must have a corresponding entry in \texttt{coordinates}. The method return a \texttt{PathwayResult} containing
\begin{itemize}
    \item the complex pathway response;
    \item the pathway name;
    \item optional numerical diagnostics;
    \item a mapping of additional observable values.
\end{itemize}

If $S_p(\omega)$ denotes the response of pathway $p$, then the returned value is the complex quantity 
\begin{equation}
    S_p(\omega)= S_p^{Re}(\omega) + iS_p^{Im}(\omega).
\end{equation}

\subsubsection{Several observables from one pathway propagation}

The method 
\begin{verbatim}
calc_pathway_observables(
    pathway,
    protocol,
    coordinates,
    observables=None
)
\end{verbatim}

propagates a pathway once and evaluates several readouts from its propagated response. The pathway's default detection is always included; additional names or \texttt{ObservableSpec} objects are normalized and deduplicated. Ordinary operator and jump-rate readouts reuse the final response. Action-detected and integrated-jump specifications can require extra post-pathway interaction branches or propagation, as described below.



\subsubsection{Component calculation}

The method \texttt{calc\_component(...)} sums selected pathways belonging to one component. If $\mathcal{P}_C$ denotes that set, the returned response is
\begin{equation}
    S_C(\omega)=\sum_{p\in\mathcal{P}_C}S_p(\omega).
\end{equation}

If \texttt{pathways=None}, all pathways belonging to the requested component are used; otherwise only the supplied subset is included.

\subsubsection{Multidimensional spectrum generation}

The method 
\begin{verbatim}
generate_spectrum(
    protocol,
    axes,
    fixed_coordinates=None,
    pathways=None,
    observables=None,
)
\end{verbatim}
evaluates the selected pathways over a multidimensional coordinate grid. The argument \texttt{axes} is a non-empty ordered mapping
\newline

\texttt{coordinate\_name: one\_dimensional\_values}
\newline

where each coordinate name must correspond to a non-identity interval in the protocol. Coordinates not included in axes must be supplied through \texttt{fixed\_coordinates}.
\newline
For example, a two-dimensional spectrum can be requested using

\begin{verbatim}
axes = {
    "omega1": omega1_values,
    "omega3": omega3_values,
}
fixed_coordinates = {
    "t2": waiting_time,
}
\end{verbatim}
The resulting array shape follows the insertion order of \texttt{axes}. The returned \texttt{SpectrumResult} contains
\begin{itemize}
    \item \texttt{axis\_names}, \texttt{axis\_values}, and \texttt{fixed\_coordinates};
    \item \texttt{pathways[name]} and component sums in \texttt{components[name]};
    \item per-pathway arrays in \texttt{observables}, indexed first by observable and then by pathway;
    \item component sums in \texttt{observable\_components}, indexed first by observable and then by component;
    \item pathway metadata and numerical diagnostics.
\end{itemize}
The default detection of each pathway remains in the legacy \texttt{pathways} and \texttt{components} mappings. Additional requested observables are stored separately, preventing one detector from overwriting another.
\newline
For every grid point and selected pathway, the solver computes $S_p(\omega)$ and stores the result independently. The component arrays are then formed by summing the corresponding pathway arrays:
\begin{equation}
    S_C(\omega) = \sum_{p \in C} S_p(\omega).
\end{equation}



\subsubsection{Coherence-order-resolved-spectrum}

The method


\begin{verbatim}
generate_nq_spectrum(
    order,
    protocol,
    axes,
    fixed_coordinates=None,
    pathways=None,
    observables=None
)
\end{verbatim}

generates a spectrum while separating pathways according to their $+NQ$ and $-NQ$ coherence-order contributions. The requested order must correspond to exactly one protocol interval whose declared coherence order has absolute value $N$. For $N>0$, the returned components include
\begin{equation}
+NQ,\qquad -NQ,\qquad NQ=(+NQ)+(-NQ).
\end{equation}

For \(N=0\), the total contribution is stored as

\begin{equation}
0Q.
\end{equation}

The same coherence-order separation is applied to the observable-component mapping. Individual pathway and observable arrays remain available in the returned result.

\subsubsection{Result containers}

\texttt{PathwayResult} stores a scalar complex \texttt{value}, the corresponding pathway, diagnostics, and observable values. \texttt{SpectrumResult} stores the multidimensional arrays described above and exposes their common \texttt{shape}. Plotting methods return a \texttt{PlotResult} containing the figure, axes, panel names, selected data, and effective configuration. These containers separate numerical results from rendering and make all aggregation steps inspectable.


\subsubsection{Jump-channel access}

The method \texttt{jump\_channel\_names()} returns the GKSL channel names supplied by the model. These names are used by instantaneous or integrated jump observables.

\subsubsection{Decay-rate diagnostic}

The rate of each channel is the Lindblad coefficient $\gamma_c$; its relation to measured decay times depends on the normalization of $\hat{L}_c$. The method
\begin{verbatim}
rates = solver.decay_rates(modes=False, max_mode_dimension=40)
\end{verbatim}
diagonalizes the Hamiltonian and reports, in its eigenbasis $\{\ket{x}\}$, the rates implied by the declared channels, read from the diagonal of the generator in the basis of dyads $\ket{x}\bra{y}$:
\begin{align}
    \Gamma_x&=\sum_c\gamma_c\left(\bra{x}\hat{L}_c^\dagger\hat{L}_c\ket{x}-|\bra{x}\hat{L}_c\ket{x}|^2\right),
    \qquad
    W_{y\leftarrow x}=\sum_c\gamma_c|\bra{y}\hat{L}_c\ket{x}|^2,\\
    \gamma_{xy}&=\sum_c\gamma_c\left[\tfrac{1}{2}\bra{x}\hat{L}_c^\dagger\hat{L}_c\ket{x}+\tfrac{1}{2}\bra{y}\hat{L}_c^\dagger\hat{L}_c\ket{y}-\operatorname{Re}\left(\bra{x}\hat{L}_c\ket{x}\bra{y}\hat{L}_c\ket{y}^*\right)\right].
\end{align}
The returned \texttt{DecayRates} object holds the energies, the dominant sector of each eigenstate, the population decay rates $\Gamma_x$ (\texttt{population\_decay}), the transfer rates $W_{y\leftarrow x}$ (\texttt{transition\_rates}), and the coherence decay rates and frequencies $\gamma_{xy}$, $\omega_{xy}$ (\texttt{coherence\_decay}, \texttt{coherence\_frequency}). The method \texttt{coherences(between=(sector\_a, sector\_b))} lists the coherences between two manifolds, for example the optical coherences between the ground and one-excitation manifolds. These secular rates are the pole half-widths of frequency-resolved intervals when the generator does not couple different dyads. With \texttt{modes=True}, the exact eigenvalues $\lambda_\nu=-\kappa_\nu-i\Omega_\nu$ of the generator are added for $D\leq$ \texttt{max\_mode\_dimension}. For the open two-level system of the Quick start, \texttt{decay\_rates()} returns $\gamma_1$ for the excited-state population and $\Gamma_2=\gamma_1/2+\gamma_\phi$ for the optical coherence; a $\hat{\sigma}_z$ channel given the rate $\gamma_\phi$ instead of $\gamma_\phi/2$ would appear as $\gamma_1/2+2\gamma_\phi$. If no dissipative channels are present, the returned tuple is empty.

\subsubsection{Direct Hilbert-space resolvent}

The public method \texttt{solve\_hilbert\_resolvent(...)} delegates the shifted Hilbert-space solve to the active backend. It is available only when the backend advertises this capability; otherwise the solver raises a capability error. The dense backend uses a direct solve and the sparse backend uses GMRES with convergence diagnostics.


\subsubsection{Memory estimation}

The method \texttt{estimate\_memory()} returns a structural memory estimate for the active backend. The result includes the total Hilbert-space dimension and the memory required for one complex state vector. For the dense Liouville backend, it additionally reports the Liouville-space dimension and the estimated memory for one density vector and one dense Liouville matrix:
\begin{equation}
    D_L = D_H^2
\end{equation}
where $D_H$ is the total Hilbert-space dimension. For the sparse backend, the estimate is structural rather than an exact prediction of the final allocation.


\subsubsection{Numerical backend architecture}
\label{app:numerical-backends}

The numerical backend is responsible for assembling the effective generator
and propagating the pathway response in the time or frequency domain.
The public \texttt{SpectroscopySolver} orchestrates the calculation, while
the backend determines how the Hilbert- and Liouville-space operations are
represented numerically.

Let

\begin{equation}
D_{\mathrm{H}}
=
\sum_{s\in\mathcal{S}}d_s
\end{equation}

denote the total Hilbert-space dimension. The corresponding vectorized
Liouville space has dimension

\begin{equation}
D_{\mathrm{L}}
=
D_{\mathrm{H}}^2.
\end{equation}

\subsubsection*{Common generator}

Both released backends use the same effective generator,

\begin{equation}
\mathcal{L}[\hat{\rho}]
=
-i[\hat{H},\hat{\rho}]
+
\sum_c
\gamma_c
\left(
\hat{L}_c\hat{\rho} \hat{L}_c^\dagger
-\frac{1}{2}
\left\{
\hat{L}_c^\dagger \hat{L}_c,\hat{\rho}
\right\}
\right).
\end{equation}

The Hamiltonian and collapse operators are assembled from the sector-resolved
blocks supplied by the external \texttt{SectorModel}. Inter-sector blocks are
allowed. The backend therefore does not assume that the global Hamiltonian is
block diagonal.

The common internal generator is represented as a matrix-free linear action
whenever possible:

\begin{equation}
\operatorname{vec}(\mathcal{L}[\hat{\rho}])
=
\mathcal{L}\operatorname{vec}(\hat{\rho}),
\end{equation}

where the matrix of \(\mathcal{L}\) acting on the vectorized density
operator carries the same symbol. The released
backends use the convention

\begin{equation}
\mathcal{L}=-i[\hat{H},\,\cdot\,]+\sum_c\gamma_c\mathcal{D}_c.
\end{equation}

A frequency-domain interval with frequency \(\omega\) and broadening
\(\eta\) solves the shifted Liouville equation

\begin{equation}
\left[
(\eta-i\omega)\mathbb{1}
-
\mathcal{L}
\right]x
=
\operatorname{vec}(\hat{\rho}).
\end{equation}

The parameter \(\eta\) is a numerical or phenomenological homogeneous
broadening and is distinct from the detection efficiency used by jump
observables.

\subsubsection*{DenseLiouvilleBackend}

\texttt{DenseLiouvilleBackend} is the reference backend for small Hilbert
spaces. During construction, it materializes the global Hamiltonian and the
complete Liouville generator:

\begin{equation}
\hat{H}\in\mathbb{C}^{D_{\mathrm{H}}\times D_{\mathrm{H}}},
\qquad
\mathcal{L}
\in
\mathbb{C}^{D_{\mathrm{L}}\times D_{\mathrm{L}}}.
\end{equation}

The time-domain propagator is explicitly constructed as

\begin{equation}
\mathcal{U}(t)
=
\exp(\mathcal{L}t),
\end{equation}

using a dense matrix exponential. A frequency-domain interval uses the
dense inverse

\begin{equation}
\mathcal{R}_\eta(\omega)
=
\left[
(\eta-i\omega)\mathbb{1}
-
\mathcal{L}
\right]^{-1}.
\end{equation}

The dense backend also provides a Hilbert-space resolvent,

\begin{equation}
R_{\mathrm{H}}(\omega,\eta)
=
\left[
(\omega+i\eta)\mathbb{1}
-
\hat{H}
\right]^{-1},
\end{equation}

which is solved using a dense linear-system solver. This Hilbert-space resolvent follows the retarded convention; for the Hamiltonian evolution of a ket it is related to the Liouville-space resolvent of the main text by $\mathcal{R}_\eta=i\,R_{\mathrm{H}}$.

Dense propagation is useful for:

\begin{itemize}
    \item small two-level or few-level systems;
    \item validating another backend;
    \item checking pathway and phase conventions;
    \item comparing against explicitly constructed Liouville matrices.
\end{itemize}

Its main limitation is memory scaling. A dense Liouville matrix contains
\(D_{\mathrm{L}}^2=D_{\mathrm{H}}^4\) complex entries. It therefore becomes
impractical rapidly as the Hilbert-space dimension increases.

The dense backend optionally caches time propagators and resolvents. Cache
sizes are bounded by the \texttt{max\_cache\_entries} option.

\subsubsection*{SparseSectorBackend}

\texttt{SparseSectorBackend} avoids constructing a global dense Hamiltonian or
a dense \(D_{\mathrm{L}}\times D_{\mathrm{L}}\) Liouville matrix. Instead, it
keeps the sector-resolved operators as linear actions.

For a Hilbert-space vector \(v\), the Hamiltonian action is assembled from the
sector blocks,

\begin{equation}
(Hv)_t
=
\sum_s \hat{H}_{ts}v_s,
\end{equation}

where \(\hat{H}_{ts}\) is the block mapping source sector \(s\) to target sector
\(t\). This allows inter-sector couplings while preserving the block
structure supplied by the model.

For time propagation, the backend uses Krylov-based actions of the form

\begin{equation}
v(t)
=
\exp(-i\hat{H}t)v(0).
\end{equation}

For a pure initial state and unitary time-domain pathway, the ket and bra
branches can be propagated separately. The response is then represented as a
rank-one dyad rather than as a full \(D_{\mathrm{L}}\)-component density
vector.

This representation is exact for a unitary evolution of a pure initial state,
up to the numerical accuracy of the Krylov propagation.

The backend switches to a Liouville-space representation when at least one of
the following conditions holds:

\begin{itemize}
    \item the protocol contains a frequency-domain interval;
    \item the generator is dissipative;
    \item the initial state is mixed.
\end{itemize}

In this case, the density response is stored as a vector of size
\(D_{\mathrm{L}}\), but the Liouville generator is still applied through a
matrix-free action. Time propagation uses a Krylov exponential action,

\begin{equation}
\operatorname{vec}(\hat{\rho}(t))
=
\exp(\mathcal{L}t)
\operatorname{vec}(\hat{\rho}(0)),
\end{equation}

without constructing the matrix of \(\mathcal{L}\) explicitly.

A frequency-domain interval is evaluated by solving

\begin{equation}
\left[
(\eta-i\omega)\mathbb{1}
-
\mathcal{L}
\right]x
=
\operatorname{vec}(\hat{\rho})
\end{equation}

with the GMRES iterative solver. The convergence residual is stored in the
pathway diagnostics. If GMRES does not converge within the requested
tolerance, the backend raises a convergence error.

The sparse backend also exposes a Hilbert-space resolvent,

\begin{equation}
\left[
(\omega+i\eta)\mathbb{1}
-
\hat{H}
\right]x
=
v,
\end{equation}

which is solved iteratively with GMRES.

The sparse backend therefore avoids the largest dense allocation, but it does
not eliminate the \(D_{\mathrm{H}}^2\) Liouville vector required for
frequency-domain, mixed-state, or dissipative calculations.

\paragraph{Options of the sparse backend.}
The options below are passed to \texttt{SpectroscopySolver} together with
\texttt{backend="sparse\_sector"} and the broadening \texttt{eta}.

\begin{itemize}
    \item \texttt{krylov\_tolerance} (default \(10^{-10}\)): relative residual
    required from each GMRES solve;
    \item \texttt{krylov\_maxiter} (default \texttt{None}): maximum number of
    GMRES iterations; \texttt{None} leaves the choice to SciPy;
    \item \texttt{resolvent\_restart} (default \texttt{None}): restart length
    of GMRES;
    \item \texttt{preconditioner} (default \texttt{None}): \texttt{None} or
    \texttt{"diagonal"}, see below; any other value raises a
    \texttt{ValueError};
    \item \texttt{cache\_propagations} (default \texttt{True}) and
    \texttt{max\_cache\_entries} (default 32): caching of propagated states and
    its bound.
\end{itemize}

\paragraph{Diagonal preconditioner.}
With \texttt{preconditioner="diagonal"}, each GMRES solve of the Liouville
resolvent is preconditioned by the inverse of the shifted diagonal of the
generator,

\begin{equation}
M_\omega^{-1}
=
\operatorname{diag}\!\left[
(\eta-i\omega)\mathbb{1}-\mathcal{L}
\right]^{-1},
\end{equation}

and each solve of the Hilbert-space resolvent by the inverse of the shifted
diagonal of \(\hat{H}\). The diagonal of \(\mathcal{L}\) is obtained from the
sector blocks without forming the matrix. The model adapters express the
operators in the eigenbasis of the Hamiltonian (sector by sector for
\texttt{ExcitationSectorModel}). The generator of a closed model, one without
collapse channels, is then diagonal, and the preconditioner is exact: every
GMRES solve converges in one iteration, for any \(\eta\) and any dimension.
With collapse channels that couple coherences the diagonal is only an
approximation: the number of iterations decreases but still depends on
\(\eta\). The preconditioner does not change the result beyond the solver
tolerance.

\begin{verbatim}
solver = SpectroscopySolver(
    backend="sparse_sector",
    eta=0.03,
    krylov_tolerance=1e-11,
    preconditioner="diagonal",
)
solver.feed_model(model, context=thermal_context)
\end{verbatim}

For the thermal spin chain of Example~2 (\(D=57\), \(\eta=0.03\)~meV), the mean
number of GMRES iterations per solve falls from about 640 to 1, and the time
for a \(3\times3\) subgrid from 58~s to 0.4~s. The measurements as a function of
\(D\) and \(\eta\) are given in Sec.~6 of the article and in
\texttt{validation/benchmarks/example2\_benchmark.ipynb}.



Both released backends are intended to represent the same physical generator.
The difference is numerical representation and linear-algebra strategy, not a
different physical model. The dense backend is useful as a reference
implementation, while the sparse backend is the preferred default for
sector-resolved calculations.

\subsubsection*{Backend selection}

The solver accepts the following backend names:

\begin{itemize}
    \item \texttt{"dense"} or \texttt{"dense\_liouville"}:
    \texttt{DenseLiouvilleBackend};
    \item \texttt{"sparse"} or \texttt{"sparse\_sector"}:
    \texttt{SparseSectorBackend};
    \item \texttt{"auto"}: automatic selection among the released backends.
\end{itemize}

The default is

\begin{verbatim}
backend="sparse_sector"
\end{verbatim}

When \texttt{backend="auto"} is selected, the solver compares the model
requirements with the declared backend capabilities. The requirements include:

\begin{itemize}
    \item pure or mixed initial-state support;
    \item unitary or Lindblad generator;
    \item time- or frequency-domain support;
    \item exactness requirements;
    \item matrix-free requirements.
\end{itemize}

The built-in automatic selection currently considers the released sparse and
dense backends. A custom backend registered through \texttt{register\_backend}
can be selected explicitly by name.





\subsubsection{Solver Summary}

The method \texttt{summary()} returns a serializable description of the constructed calculation without exposing the underlying external model. The summary includes
\begin{itemize} 
    \item the selected backend;
    \item backend capabilities and support;
    \item model requirements and capabilities;
    \item the sector labels and dimensions;
    \item the total Hilbert-space dimension;
    \item the stationarity residual of the reference state;
    \item the active pathway metadata;
    \item the transition-operator decomposition and its \texttt{transition\_summary()} when they are advertised by the
    model.
\end{itemize}
This method is intended for inspection, reproducibility records, and debugging of backend or pathway configuration. 

\subsection{\texttt{ObservableSpec}: operator, jump, and action detection}
\label{sec:observables}

\texttt{ObservableSpec} separates the measured quantity from pathway propagation:
\begin{verbatim}
ObservableSpec(
    name,
    kind="operator",
    operator=None,
    channel=None,
    efficiency=1.0,
    time_window=None,
    n_steps=101,
    fourth_interaction=None,
    integration="exact",
)
\end{verbatim}
The efficiency must lie in $[0,1]$. Integrated observables require $0\leq t_0\leq t_f$; \texttt{integration} is \texttt{"exact"} (default) or \texttt{"trapezoid"}, and \texttt{n\_steps}, at least two, is used only by the trapezoidal option.

\subsubsection{Named operator observables}

For \texttt{kind="operator"}, \texttt{operator} identifies a named model-side observable $\hat{O}$. If it is omitted, the observable name is used. For the final pathway response $\hat{\rho}_p$,
\begin{equation}
    S_{O,p}=c_p\operatorname{Tr}(\hat{O}\hat{\rho}_p),
\end{equation}
where $c_p$ is the pathway response prefactor defined in Section~\ref{sec:pathway-conventions}. Multiple named operators may be exposed by either provided model adapter.

\subsubsection{Instantaneous and integrated jump observables}

For \texttt{kind="jump\_rate"}, \texttt{channel} identifies a GKSL channel. The pathway-resolved detected rate is
\begin{equation}
    I_{c,p}(t)=c_p\varepsilon_{\mathrm{det}}\gamma_c
    \operatorname{Tr}\!\left(\hat{L}_c^\dagger \hat{L}_c\hat{\rho}_p(t)\right).
\end{equation}
The channel rate $\gamma_c$ remains separate from $\hat{L}_c$. Hence the channel must be one of \texttt{solver.jump\_channel\_names()}. A pathway contribution can be complex because it includes $c_p$, even though the rate of a physical density matrix is real and non-negative.

For \texttt{kind="integrated\_jump"}, the detected rate is integrated over a detection window that starts after the final pathway interval,
\begin{equation}
    N_{c,p}=\int_{t_0}^{t_f}I_{c,p}(t)\,dt
    =c_p\,\langle\!\langle \hat{X}_{c}|\hat{\rho}_p\rangle\!\rangle,
    \qquad
    \hat{X}_{c}=\varepsilon_{\mathrm{det}}\gamma_c\int_{t_0}^{t_f}\mathcal{U}^\dagger(t)\big[\hat{L}_c^\dagger\hat{L}_c\big]\,dt.
\end{equation}
With the default \texttt{integration="exact"}, the integrated detector $\hat{X}_c$ is computed once per channel and window, without quadrature, as the last column of the exponential of the augmented generator $\big(\begin{smallmatrix}\mathcal{L}^\dagger & \vec{v}\\ 0&0\end{smallmatrix}\big)$ with $\vec{v}=\operatorname{vec}(\gamma_c\hat{L}_c^\dagger\hat{L}_c)$, followed by $\exp(\mathcal{L}^\dagger t_0)$ when $t_0>0$. The dense backend uses a dense exponential; the sparse backend applies the augmented generator matrix-free with \texttt{expm\_multiply}. The result is cached, so that each pathway state and grid point costs one inner product. With \texttt{integration="trapezoid"}, the propagated response is instead sampled at \texttt{n\_steps} uniformly spaced times and integrated with the trapezoidal rule; this legacy option has an error of order $(\Delta t)^2$ and a cost proportional to \texttt{n\_steps} per grid point. The convenience constructor
\begin{verbatim}
ObservableSpec.mean_jump(
    name,
    channel,
    *,
    time_window=(t0, tf),
    efficiency=1.0,
    n_steps=101,
    fourth_interaction=None,
    integration="exact",
)
\end{verbatim}
creates this specification. \texttt{observable.instantaneous()} returns the corresponding instantaneous jump-rate detector.

\subsubsection{Action detection through a fourth interaction}

The field \texttt{fourth\_interaction} requests one or more additional perturbative interactions after the final pathway interval and before detection. It accepts an \texttt{Interaction}, a pathway label such as \texttt{"Ku"}, a mapping, or a sequence of alternatives. The convenience constructor is
\begin{verbatim}
ObservableSpec.action(
    name,
    fourth_interaction,
    *,
    kind="operator",
    operator=None,
    channel=None,
    efficiency=1.0,
    time_window=None,
    n_steps=101,
    integration="exact",
)
\end{verbatim}
For alternatives $a\in\mathcal{I}_{\mathrm{rd}}$, the solver applies each interaction to the propagated response and sums the branches before the requested detector,
\begin{equation}
    S_{O,p}^{(4)}=
    c_p\sum_{a\in\mathcal{I}_{\mathrm{rd}}}c_{\mathrm{rd},a}
    \operatorname{Tr}\!\left[\hat{O}\,\mathcal{V}_a(\hat{\rho}_p)\right],
    \qquad
    c_{\mathrm{rd},a}=\begin{cases}+i,&\text{ket action},\\-i,&\text{bra action}.
    \end{cases}
\end{equation}
This fourth interaction adds no protocol interval or spectral axis. For an integrated jump observable, every projected branch is contracted with the same integrated detector (or, with the trapezoidal option, propagated across the detection window before integration). Dedicated properties expose the normalized projection and whether action detection is active; \texttt{without\_projection()} returns the underlying detector specification.

\subsubsection{Normalization and aliases}

\texttt{normalize\_observables(...)} accepts a string or specification, a sequence, or a mapping from aliases to specifications. It returns a validated tuple with unique names; a later duplicate replaces the earlier specification without changing its original position. The aliases \emph{observable} and \emph{trace} map to the operator kind. The aliases \emph{mean jump} and \emph{integrated rate} map to the integrated-jump kind. At least one effective observable is required. The solver retains the pathway's default detection and appends requested observables to the structured result mappings.


\subsection{\textbf{Pathway conventions}}
\label{sec:pathway-conventions}

A pathway is an ordered sequence of light--matter interactions acting on the ket or bra branch of the Liouville-space density operator. In QuDPy-FDGF, each interaction is represented by an \texttt{Interaction} object and each complete pathway by a \texttt{FrequencyPathway} object.



\subsubsection{Interaction fields}

An interaction is defined by

\begin{verbatim}
Interaction(
    label,
    operator="light_matter",
    pulse_index=None,
    source_sector=None,
    target_sector=None,
)
\end{verbatim}

The fields have the following meanings:

\begin{itemize}
    \item \texttt{label} identifies the side and transition direction;
    \item \texttt{operator} gives the named transition operator;
    \item \texttt{pulse\_index} identifies the associated pulse;
    \item \texttt{source\_sector} optionally constrains the source sector;
    \item \texttt{target\_sector} optionally constrains the target sector.
\end{itemize}

The side of an interaction is determined by its label:

\begin{equation}
\texttt{Ku},\texttt{Kd}\ \Rightarrow\ \text{ket},
\qquad
\texttt{Bu},\texttt{Bd}\ \Rightarrow\ \text{bra}.
\end{equation}

The transition direction applied internally is

\begin{equation}
\texttt{Ku},\texttt{Bu}\ \Rightarrow\ \hat{\mu}_{+},
\qquad
\texttt{Kd},\texttt{Bd}\ \Rightarrow\ \hat{\mu}_{-}.
\end{equation}

The apparent inversion for bra-side interactions results from the fact that right multiplication of a density matrix is implemented by applying the adjoint operator to the stored bra vector. This implementation preserves the coherence-order convention defined in the next subsection while allowing ket and bra branches to be propagated using ordinary Hilbert-space vectors.


\subsubsection{Coherence-order history}

For a dyad $\ket{N_{\mathrm{ket}},a}\bra{N_{\mathrm{bra}},b}$, the signed coherence order is
\begin{equation}
q=N_{\mathrm{ket}}-N_{\mathrm{bra}},
\end{equation}
where the two $N$ values are excitation-manifold indices, not counts of ket- and bra-side interactions. The coherence-order change associated with each interaction is

\begin{equation}
\Delta q(\texttt{Ku})=+1, \qquad \Delta q(\texttt{Kd})=-1, \qquad \Delta q(\texttt{Bu})=-1, \qquad \Delta q(\texttt{Bd})=+1.
\end{equation}

For an ordered interaction sequence \((I_1,I_2,\ldots,I_N)\), the coherence order after interaction \(n\) is

\begin{equation}
q_n=\sum_{m=1}^{n}\Delta q(I_m),\qquad q_0=0.
\end{equation}

The complete coherence-order history is therefore

\begin{equation}
(q_1,q_2,\ldots,q_N).
\end{equation}

This history is automatically generated from the interaction labels when \texttt{coherence\_orders} is omitted. If a coherence-order tuple is provided explicitly, it must exactly match the values implied by the interaction sequence.

The coherence-order history is used by \texttt{generate\_nq\_spectrum} to separate \(+NQ\), \(-NQ\), and total \(NQ\) contributions.




\subsubsection{FrequencyPathway fields}

A complete pathway is defined by

\begin{verbatim}
FrequencyPathway(
    name,
    interactions,
    component="custom",
    amplitude=1.0,
    prefactor=None,
    coherence_orders=(),
    detection="polarization",
)
\end{verbatim}

The number of interactions in a pathway must match the number of intervals in the associated \texttt{SpectroscopyProtocol}. Each interaction is paired with one propagation interval.

The component name is normalized to lower case. The nonrephasing aliases are internally mapped to the canonical unrephasing component. Other names may be used for custom groupings.


\subsubsection{Pathway response prefactor}

If an explicit \texttt{prefactor} is provided, it is used directly as the complex multiplier of the pathway response. Otherwise, the solver constructs the response prefactor from the pathway amplitude, the number of interactions, and the number of bra-side interactions.

Let \(N_{\mathrm{int}}\) be the total number of interactions and \(N_{\mathrm{bra}}\) the number of bra-side interactions. The default
prefactor is

\begin{equation}
c_p = a_p\, i^{N_{\mathrm{int}}} (-1)^{N_{\mathrm{bra}}},
\end{equation}

where \(a_p\) is the complex value of \texttt{amplitude}.

More generally,

\begin{equation}
c_p=
\begin{cases}
\texttt{prefactor},& \text{if an explicit prefactor is supplied},\\[4pt]
\texttt{amplitude}\, i^{N_{\mathrm{int}}} (-1)^{N_{\mathrm{bra}}},& \text{otherwise}.
\end{cases}
\end{equation}


The final pathway response and all associated observables are multiplied by
\(c_p\). Therefore, experimental field phases or alternative diagrammatic
sign conventions should be encoded through \texttt{amplitude} or an explicit
\texttt{prefactor}.

\subsubsection{UFSS diagram translation}

The helper

\begin{verbatim}
translate_ufss_diagrams(
    diagrams,
    component="custom",
    names=None,
    operator="light_matter",
    amplitudes=None,
    prefactors=None,
    detection="polarization",
)
\end{verbatim}

translates UFSS-style diagrams into QuDPy-FDGF pathways. Each diagram is represented as an ordered sequence of pairs

\begin{equation}
(\texttt{label},\texttt{pulse\_index}),
\end{equation}

for example

\begin{verbatim}
[
    ("Ku", 1),
    ("Bd", 2),
    ("Ku", 3),
]
\end{verbatim}

The helper converts every pair into an \texttt{Interaction}, preserves the
interaction order, assigns unique pathway names, and constructs the
corresponding \texttt{FrequencyPathway} objects.

This translation preserves the UFSS coherence-order convention

\begin{equation}
\Delta q = \{+1,-1,-1,+1\}
\end{equation}

for the labels
\(\{\texttt{Ku},\texttt{Kd},\texttt{Bu},\texttt{Bd}\}\), respectively.


\subsection{\texttt{SpectroscopyPlotter}}
\label{sec:plotter}

\texttt{SpectroscopyPlotter} is a rendering-only utility for objects returned by \texttt{SpectroscopySolver}. It selects, phase-rotates, normalizes, and renders already calculated data; it never changes the numerical result. Matplotlib is therefore an optional dependency of the plotting layer only.

\subsubsection{Constructor}

The constructor is

\begin{verbatim}
SpectroscopyPlotter(
    w_list=None,
    detection_phase=0.0,
)
\end{verbatim}

The optional \texttt{w\_list} provides a default frequency axis for plotting raw spectra with \texttt{plot\_contourf\_multi\_spectra} or \texttt{plot\_1d}. 

The parameter \texttt{detection\_phase} specifies a display-phase correction. For a complex signal $S$, the plotter applies
\begin{equation}
    S_{\phi}=\exp(i\phi_{\mathrm{det}})S
\end{equation}
before separating real, imaginary, and absolute views. A plotting configuration may override the constructor value locally.


\subsubsection{Selecting data from a \texttt{SpectrumResult}}

The method
\begin{verbatim}
select_spectrum_result_data(
    result,
    source="pathways",
    names=None,
    pathways=None,
    totals="auto",
    observable=None,
)
\end{verbatim}

selects panels from a two-dimensional result. The source may contain pathways, components, observables, or observable components; a user mapping is also accepted. For observable sources, \texttt{observable} selects the detector. \texttt{names} or the legacy \texttt{pathways} argument restricts the panels, and \texttt{totals} controls component-total panels. The first result axis is plotted vertically and the second horizontally.



\subsubsection{Dictionary-based plotting configuration}

The principal entry point is
\begin{verbatim}
plot_spectrum_result(result, params=None, **overrides)
\end{verbatim}
where keyword overrides take precedence over \texttt{params}. Important options include the data \texttt{source}, selected names, \texttt{observable}, view, normalization, detection phase, diagonal guides, quadrant, colormap, contour levels, labels, title, saving, and display. The method returns a \texttt{PlotResult} rather than only a Matplotlib figure.

\subsubsection{Views, normalization, and guides}

The \texttt{view} option accepts \texttt{"real"}, \texttt{"imag"}, \texttt{"abs"}, or \texttt{"all"}. Normalization may be independent per panel, shared per row, global, or disabled through \texttt{"panel"}, \texttt{"row"}, \texttt{"global"}, and \texttt{"none"}; legacy aliases are normalized internally. Real and imaginary panels use symmetric color limits around zero, whereas absolute-value panels use non-negative limits.

Diagonal reference lines may be selected automatically or requested as rephasing, unrephasing, or both. \texttt{plot\_quadrant} accepts the full plane, a zoomed region, or quadrants 1--4. Convenience wrappers cover pathway panels, real/imaginary/absolute views, result contours, and multiple coherence orders. Raw-array utilities remain available for contour and one-dimensional plots.

\end{document}
